% Author: Suwon Lee, Kookmin University
% VFG 경로추종 시스템 배경이론 보고서
\documentclass[11pt,a4paper]{article}

\usepackage{fontspec}
\usepackage{amsmath,amssymb,mathtools}
\usepackage{graphicx}
\usepackage{booktabs}
\usepackage[a4paper,margin=2.5cm]{geometry}
\usepackage{hyperref}
\usepackage{enumitem}
\usepackage{caption}

% Korean font support (XeLaTeX)
\usepackage{kotex}

% Math commands
\newcommand{\vect}[1]{\mathbf{#1}}
\newcommand{\mat}[1]{\mathbf{#1}}
\DeclareMathOperator{\sgn}{sgn}
\DeclareMathOperator{\atantwo}{atan2}

\title{VFG 경로추종 시스템 --- 배경이론 보고서\\
\large Vector Field Guidance Path Following: Background Theory}
\author{이수원 \\ 국민대학교 기계공학부}
\date{2026년 3월}

\begin{document}
\maketitle

\begin{abstract}
이 보고서는 LIMO 모바일 로봇을 위한 VFG(Vector Field Guidance) 기반 경로추종 시스템의 배경이론을 학부생 수준에서 설명한다.
경로추종 문제 정의, VFG 유도 원리, PID 및 LPV $\mathcal{H}_\infty$ 제어기 개념, 시뮬레이션 결과 비교를 다룬다.
$\mathcal{H}_\infty$ 합성의 수학적 세부사항은 생략하고, 학생이 시뮬레이션 도구를 활용하여 경로추종 실험을 수행하는 데 필요한 직관적 이해를 제공하는 것을 목표로 한다.
\end{abstract}

\tableofcontents
\newpage

% Author: Suwon Lee, Kookmin University
\section{경로추종 문제 정의}
\label{sec:problem}

\subsection{경로(Path)란?}

경로추종(path following)에서 \textbf{경로}란 로봇이 따라가야 할 2차원 기하학적 곡선이다.
시간에 대한 정보는 포함하지 않으며, 호 길이(arc-length) $s$로 매개변수화한다:
\begin{equation}
    \vect{P}(s) = \begin{bmatrix} x(s) \\ y(s) \end{bmatrix}, \quad s \in [0, L]
\end{equation}
여기서 $L$은 경로의 전체 길이이다.

경로 위의 각 점에서 다음 기하 정보를 정의할 수 있다:
\begin{itemize}[nosep]
    \item \textbf{접선 벡터} $\vect{t}(s)$: 경로 진행 방향의 단위 벡터
    \item \textbf{법선 벡터} $\vect{n}(s)$: 접선에 수직, 왼쪽 방향 (반시계)
    \item \textbf{곡률} $\kappa(s)$: 경로의 ``휘어진 정도''. $\kappa = 1/R$ (원호의 경우 반지름의 역수)
    \item \textbf{경로 방향각} $\psi_\text{path}(s) = \atantwo(y'(s), x'(s))$
\end{itemize}

\subsection{추종 오차}

로봇 위치 $\vect{q} = (X, Y)$에서 경로 위 가장 가까운 점 $\vect{P}(s^*)$를 찾으면,
두 가지 오차를 정의할 수 있다:

\begin{itemize}[nosep]
    \item \textbf{횡방향 오차} (cross-track error): $e_d = (\vect{q} - \vect{P}(s^*)) \cdot \vect{n}(s^*)$
    \item \textbf{방향 오차} (heading error): $e_\psi = \psi_\text{des} - \psi$
\end{itemize}

$e_d > 0$이면 로봇이 경로 왼쪽에, $e_d < 0$이면 오른쪽에 있다.
$e_\psi > 0$이면 로봇이 원하는 방향보다 시계 방향으로 벗어나 있다.

경로추종의 목표는 두 오차를 모두 0으로 수렴시키는 것이다:
\begin{equation}
    e_d \to 0, \quad e_\psi \to 0
\end{equation}

\subsection{자전거 모델 (Bicycle Model)}

LIMO 로봇은 Ackermann 조향 구조를 갖는 1:10 스케일카이다.
저속에서는 운동학적 자전거 모델(kinematic bicycle model)로 충분히 근사할 수 있다.

\begin{equation}
\begin{aligned}
    \dot{X} &= v \cos(\psi + \beta) \\
    \dot{Y} &= v \sin(\psi + \beta) \\
    \dot{\psi} &= \frac{v}{L} \cos\beta \cdot \tan\delta \\
    \dot{\delta} &= \frac{\delta_\text{cmd} - \delta}{\tau_\delta}
\end{aligned}
\label{eq:bicycle}
\end{equation}
여기서 $v$: 종방향 속도, $\psi$: 방향각, $\delta$: 조향각,
$L = 0.2$\,m: 축거, $\tau_\delta = 0.07$\,s: 조향 시정수, $\beta$: 슬립각이다.

핵심 매개변수:
\begin{itemize}[nosep]
    \item 축거 $L = 0.2$\,m --- 앞바퀴와 뒷바퀴 축 사이 거리
    \item 조향 시정수 $\tau_\delta = 0.07$\,s --- 조향 명령이 실제 조향각에 반영되는 지연
    \item 최대 조향각 $\delta_\text{max} = 0.5$\,rad ($\approx 28.6°$)
\end{itemize}

조향 명령 $\delta_\text{cmd}$를 입력하면, 1차 지연을 거쳐 실제 조향각 $\delta$에 반영된다.
이 조향각이 차량의 방향을 변화시키고, 결과적으로 경로추종 오차를 줄인다.

% Author: Suwon Lee, Kookmin University
\section{VFG 유도 원리}
\label{sec:vfg}

\subsection{Vector Field Guidance란?}

VFG(Vector Field Guidance)는 경로 주변에 ``흐름의 장''을 형성하여
로봇이 자연스럽게 경로로 수렴하고 경로를 따라가도록 유도하는 기법이다.

직관적으로, 강물의 흐름을 상상하면 된다:
\begin{itemize}[nosep]
    \item 경로에서 멀리 있으면 → 강하게 경로 쪽으로 당기는 흐름 (수렴)
    \item 경로 위에 있으면 → 경로 방향으로 흘러가는 흐름 (추종)
\end{itemize}

이 두 성분의 자연스러운 전환이 VFG의 핵심이다.

\subsection{수렴/추종 벡터 분해}

경로 위 가장 가까운 점 $\vect{P}(s^*)$에서 두 방향 벡터를 정의한다:
\begin{itemize}[nosep]
    \item \textbf{추종 벡터} $\vect{v}_\text{trav} = \vect{t}(s^*)$: 경로 접선 방향
    \item \textbf{수렴 벡터} $\vect{v}_\text{conv} = -\sgn(e_d) \cdot \vect{n}(s^*)$: 경로를 향하는 방향
\end{itemize}

이 두 벡터를 거리에 따라 가중 합산한다:
\begin{equation}
    \vect{v}_\text{des} = k_\text{conv}(d) \cdot \vect{v}_\text{conv} + k_\text{trav}(d) \cdot \vect{v}_\text{trav}
    \label{eq:vfg}
\end{equation}
여기서 $d = |e_d|$는 경로까지의 거리이다.

\subsection{tanh 가중 함수}

가중 계수는 $\tanh$ 함수로 정의한다:
\begin{equation}
    k_\text{conv}(d) = \tanh(k_e \cdot d), \quad
    k_\text{trav}(d) = \sqrt{1 - k_\text{conv}^2(d)}
    \label{eq:weights}
\end{equation}

$k_e$는 \textbf{수렴 이득}(convergence gain)으로, 경로로 수렴하는 강도를 결정한다.

\begin{itemize}[nosep]
    \item $d \gg 0$: $\tanh(k_e d) \approx 1$ → 수렴 우세 (경로 쪽으로 이동)
    \item $d \approx 0$: $\tanh(k_e d) \approx 0$ → 추종 우세 (경로를 따라 이동)
    \item $k_e$ 증가: 더 빠르게 수렴, 그러나 너무 크면 오버슈트 발생
\end{itemize}

기본값 $k_e = 3.0$은 대부분의 상황에서 적절한 수렴 성능을 보인다.

\subsection{원하는 방향각}

VFG의 출력은 원하는 방향각(desired heading)이다:
\begin{equation}
    \psi_\text{des} = \atantwo\!\left(v_{\text{des},y},\; v_{\text{des},x}\right)
\end{equation}

이 $\psi_\text{des}$와 실제 로봇 방향각 $\psi$의 차이가 방향 오차 $e_\psi$이며,
제어기가 이 오차를 줄이도록 조향 명령을 생성한다.

\subsection{VFG의 역할 구분}

VFG는 ``어디로 향해야 하는가''를 결정하고 (유도, guidance),
제어기는 ``조향을 얼마나 할 것인가''를 결정한다 (제어, control).

이 분리 구조 덕분에:
\begin{itemize}[nosep]
    \item VFG는 경로 기하에만 의존 → 제어기 교체 시에도 유도 모듈 재사용 가능
    \item 제어기는 방향 오차만 처리 → 다양한 제어 기법 적용 가능
    \item 공정한 비교: 같은 VFG 위에서 PID와 LPV $\mathcal{H}_\infty$를 직접 비교할 수 있음
\end{itemize}

% Author: Suwon Lee, Kookmin University
\section{제어기 개념}
\label{sec:controller}

\subsection{PID + 곡률 피드포워드 (PID-FF)}

가장 직관적인 제어기이다.
방향 오차 $e_\psi = \psi_\text{des} - \psi$에 비례-미분 제어를 적용하고,
경로 곡률에 대한 피드포워드 항을 추가한다:

\begin{equation}
    \delta_\text{cmd} = \underbrace{K_P \cdot e_\psi + K_D \cdot \dot{e}_\psi}_{\text{PID 피드백}}
    + \underbrace{\arctan(L \cdot \kappa)}_{\text{곡률 피드포워드}}
    \label{eq:pid}
\end{equation}

\begin{itemize}[nosep]
    \item $K_P = 2.0$: 비례 이득. 오차에 비례하여 조향
    \item $K_D = 0.3$: 미분 이득. 오차 변화율에 반응하여 진동 억제
    \item 피드포워드: 경로가 휘어져 있으면 그에 맞는 조향을 미리 적용
\end{itemize}

PID-FF는 구현이 간단하고 직관적이지만, 속도와 곡률이 변하는 상황에서
이득이 고정되어 있어 최적 성능을 보장하지 못한다.

\subsection{LPV $\mathcal{H}_\infty$ 제어기 (블랙박스 관점)}

\subsubsection{$\mathcal{H}_\infty$ 제어란?}

$\mathcal{H}_\infty$ 제어는 ``최악의 경우에도 성능을 보장''하는 강인 제어 기법이다.
일상적인 비유: \textbf{강인한 자동조종 장치}.

\begin{itemize}[nosep]
    \item 센서 잡음이 있어도 안정적으로 동작
    \item 모델과 실제 로봇 사이 차이가 있어도 강인
    \item 주파수 영역에서 성능/강인성 트레이드오프를 명시적으로 설계
\end{itemize}

$\mathcal{H}_\infty$ 합성의 수학적 세부 사항은 이 보고서의 범위를 벗어난다.
학생은 사전 계산된 제어기(JSON 파일)를 그대로 사용하면 된다.

\subsubsection{LPV 스케줄링}

LPV(Linear Parameter-Varying)는 ``상황에 따라 이득이 자동으로 변하는''
제어기이다.

스케줄링 매개변수 $\rho$를 다음과 같이 정의한다:
\begin{equation}
    \rho = |\kappa| \cdot v
\end{equation}

$\rho$의 물리적 의미:
\begin{itemize}[nosep]
    \item $\rho = 0$: 직선 구간 또는 정지 → 제어기가 보수적으로 동작
    \item $\rho$ 큼: 고속 + 고곡률 → 제어기가 적극적으로 반응
\end{itemize}

6개의 꼭짓점(vertex) 제어기 ($\rho = 0, 1, 2, 3, 4, 5$)가 사전 설계되어 있으며,
실시간에서 현재 $\rho$ 값에 따라 인접 꼭짓점 사이를 선형 보간한다.

\subsubsection{입출력 구조}

학생이 알아야 할 것은 입출력 관계뿐이다:

\begin{itemize}[nosep]
    \item \textbf{입력}: 방향 오차 $e_\psi$, 측정 조향각 $\delta_\text{meas}$, 스케줄링 변수 $\rho$
    \item \textbf{출력}: 조향 명령 $\delta_\text{cmd}$
    \item \textbf{내부 상태}: 6차 상태공간 모델 (자동 관리)
\end{itemize}

부호 규칙: 학생은 $e_\psi = \psi_\text{des} - \psi$를 그대로 넘기면 된다.
내부적으로 MATLAB 관례에 맞게 부호가 자동 변환된다.

\subsection{두 제어기의 비교}

\begin{table}[h]
    \centering
    \caption{PID-FF와 LPV $\mathcal{H}_\infty$ 비교}
    \begin{tabular}{lcc}
        \toprule
        항목 & PID-FF & LPV $\mathcal{H}_\infty$ \\
        \midrule
        구현 난이도 & 쉬움 & 블랙박스 (JSON 로드) \\
        이득 조정 & 수동 ($K_P, K_D$) & 자동 ($\rho$ 기반) \\
        잡음 강인성 & 보통 & 높음 (설계에 반영) \\
        고속 성능 & 고정 이득 한계 & $\rho$ 스케줄링으로 적응 \\
        저속 직선 & 우수 & 유사 \\
        파라미터 수 & 2 ($K_P, K_D$) & 0 (사전 설계) \\
        \bottomrule
    \end{tabular}
    \label{tab:comparison}
\end{table}

핵심 메시지: 저속 직선 구간에서는 PID로 충분하지만,
고속이나 급곡선에서는 LPV $\mathcal{H}_\infty$가 체계적인 이점을 보인다.
다음 장의 시뮬레이션 결과에서 이를 정량적으로 확인한다.

% Author: Suwon Lee, Kookmin University
\section{시뮬레이션 결과 비교}
\label{sec:results}

\subsection{실험 설정}

시뮬레이션은 \texttt{vfg\_pathfollowing} 패키지를 사용하여 수행한다.
동일한 VFG 유도 모듈($k_e = 3.0$) 위에서 PID-FF와 LPV $\mathcal{H}_\infty$를 비교한다.

경로: 직선 $\to$ 원호(R=0.5\,m, $\theta = 90°$) $\to$ 직선 (step curvature path)

이 경로는 갑작스러운 곡률 변화를 포함하여, 제어기의 과도 응답 성능을 평가하기에 적합하다.

\subsection{속도별 성능}

표~\ref{tab:results}에 속도별 RMS 방향 오차와 최대 방향 오차를 정리하였다.

\begin{table}[h]
    \centering
    \caption{속도별 방향 오차 비교 (Step curvature path, R=0.5\,m)}
    \begin{tabular}{ccccc}
        \toprule
        & \multicolumn{2}{c}{RMS $e_\psi$ [deg]} & \multicolumn{2}{c}{Max $e_\psi$ [deg]} \\
        \cmidrule(lr){2-3} \cmidrule(lr){4-5}
        $v$ [m/s] & PID-FF & LPV $\mathcal{H}_\infty$ & PID-FF & LPV $\mathcal{H}_\infty$ \\
        \midrule
        1.0 & 4.17 & 2.83 & 15.8 & 14.2 \\
        1.5 & 3.77 & 2.34 & 18.1 & 11.5 \\
        2.0 & 3.40 & 2.16 & 19.4 & 8.7 \\
        2.5 & 3.39 & 2.14 & 20.2 & 7.3 \\
        \bottomrule
    \end{tabular}
    \label{tab:results}
    \vspace{0.5em}
    \small{※ 200회 Monte Carlo 시뮬레이션 (센서 잡음 포함) 평균값}
\end{table}

\subsection{주요 관찰}

\begin{enumerate}[nosep]
    \item \textbf{저속 ($v = 1.0$\,m/s)}: 두 제어기 모두 양호한 성능. PID와 LPV의 차이가 작다.
    이 영역에서 $\rho = |\kappa| \cdot v$가 작아 LPV 스케줄링의 이점이 제한적이다.

    \item \textbf{중속 ($v = 1.5$\,m/s)}: LPV가 RMS에서 38\% 우위를 보이기 시작한다.
    곡률 구간 진입 시 $\rho$가 증가하면서 제어기 대역폭이 자동으로 확대된다.

    \item \textbf{고속 ($v \geq 2.0$\,m/s)}: LPV의 우위가 뚜렷하다.
    특히 \textbf{최대 오차}에서 큰 차이가 난다 --- $v = 2.0$에서 PID 19.4° vs LPV 8.7° (55\% 감소).
    PID는 고정 이득이므로 속도가 증가해도 동일한 이득을 사용하지만,
    LPV는 $\rho$ 증가에 따라 대역폭을 자동으로 높여 더 빠르게 반응한다.

    \item \textbf{센서 잡음 영향}: 두 제어기 모두 잡음($\sigma_\psi = 0.03$\,rad) 하에서 안정적이다.
    LPV는 $\mathcal{H}_\infty$ 설계 시 잡음 모델이 포함되어 있어 추가적인 강인성을 보인다.
\end{enumerate}

\subsection{시뮬레이션 실행 방법}

노트북 \texttt{02\_pid\_vs\_lpvhinf.ipynb}에서 위 결과를 직접 재현할 수 있다:

\begin{verbatim}
from vfg_pathfollowing import Simulator, StepCurvaturePath

path = StepCurvaturePath(R=0.5, theta_arc=1.57)
sim = Simulator(path, speed=2.0)
results = sim.compare(T=15.0)
Simulator.plot_comparison(results, path=path)
\end{verbatim}

\subsection{정리}

VFG 유도 + LPV $\mathcal{H}_\infty$ 제어 조합은 PID-FF 대비 고속/고곡률 구간에서
체계적인 성능 향상을 보인다.
이는 $\rho$-스케줄링을 통해 주행 조건에 맞게 제어기 대역폭을 자동 조절하기 때문이다.
학생은 노트북을 통해 다양한 속도, 경로, 파라미터에서 이 차이를 직접 확인할 수 있다.


\end{document}
